\documentclass[10pt]{article}
\usepackage[utf8]{inputenc}
\usepackage[margin=1in]{geometry}
\usepackage{amsmath}
\usepackage{booktabs}
\usepackage{array}
\usepackage{enumitem}
\usepackage{hyperref}

% Reduce spacing around sections
\makeatletter
\renewcommand\section{\@startsection{section}{1}{\z@}%
  {-1.5ex \@plus -0.5ex \@minus -.2ex}%
  {0.8ex \@plus.1ex}%
  {\normalfont\Large\bfseries}}
\makeatother

\title{Quantifying CNN Robustness to Weather Corruptions in Traffic Sign Recognition}
\author{Ayah Halabi, Setara Nusratty, Mark Rahal, Hochan Son}
\date{Stats 426 | Timeline: 4 weeks}
\begin{document}
\maketitle
\section{Objective}
Deep CNNs achieve 99\% accuracy on clean traffic sign datasets but degrade significantly under real-world conditions. We quantify this robustness gap by measuring ResNet-50's performance on synthetic weather corruptions (rain, fog, low brightness) to identify which conditions and sign types are most vulnerable.
\section{Methodology}
We use GTSRB (German Traffic Signs) and LISA (US Traffic Signs) datasets, focusing on 5 classes (Stop, Yield, Speed 30/50/70 km/h) with $224 \times 224$ RGB images. The data is split into 80\% training and 20\% validation from the training set, with a separate test set reserved for corruption experiments.

We fine-tune ResNet-50 (pretrained on ImageNet) on the clean 5-class subset using Adam optimizer (lr=$10^{-4}$), 15 epochs, and early stopping, targeting $>$90\% clean validation accuracy.

Using the Albumentations library, we apply three corruptions to the test set only: (1) rain with synthetic droplets and reduced contrast, (2) fog with atmospheric scattering (fog\_coef=0.5), and (3) low brightness with 30\% reduction.

We evaluate performance using degradation score ($\text{Accuracy}_{\text{corrupted}} / \text{Accuracy}_{\text{clean}}$), $5 \times 5$ confusion matrices per corruption type, and a vulnerability index ranking per-class susceptibility.
\section{Expected Results}
\begin{table}[h]
\centering
\begin{tabular}{lll}
\toprule
\textbf{Metric} & \textbf{Expected Range} & \textbf{Hypothesis} \\
\midrule
Clean accuracy & 92--98\% & Baseline performance \\
Rain degradation & $<0.85$ & Moderate impact \\
\textbf{Fog degradation} & $\mathbf{<0.75}$ & \textbf{Most damaging} (loss of high-frequency details) \\
Brightness degradation & $<0.90$ & Least impact \\
\bottomrule
\end{tabular}
\end{table}
\textbf{Key Prediction:} Speed limit signs (text-dependent) will be more vulnerable than Stop/Yield (shape-dependent).
\section{Success Criteria}
\textbf{Minimum:} ResNet-50 trained ($>$85\%), degradation scores for 3 corruptions, confusion matrices, 6-page report
\textbf{Ideal:} $>$92\% clean accuracy, clear vulnerability patterns identified, publication-quality analysis
\section{References}
\begin{enumerate}[nosep]
    \item Stallkamp et al. (2012). Man vs. computer: Benchmarking ML for traffic sign recognition. \textit{Neural Networks}, 32, 323--332.
    \item Hendrycks \& Dietterich (2019). Benchmarking neural network robustness to common corruptions. \textit{ICLR}.
    \item Qiao, X. (2023). Traffic sign recognition based on CNN. \textit{Procedia Computer Science}, 220, 107--114.
\end{enumerate}

\end{document}
